\subsection{有限 primitive 手术的主极点不变性与 Abel 有限部变换}\label{subsec:conclusion-finite-primitive-surgery-pole-abel-law}
沿用命题 \ref{prop:primitive-orbit-surgery}、命题 \ref{prop:atomic-2-orbit-split-logM} 以及真实输入加权 \(\zeta\)-函数的记号。前述链条已经把单个 primitive 原子因子的剥离在 \(\zeta\)、primitive 与 Abel 有限部三层上完全显式化。由此，对有限个原子 Euler 因子的反复手术还可进一步压缩为一条同时控制主极点位置、主留数与 Abel 有限部的统一变换律。

\begin{theorem}[有限 primitive 手术的主极点不变性与 Abel 有限部变换]\label{thm:conclusion-finite-primitive-surgery-pole-finitepart-transform}
设在某个参数域上有分解
\[
\zeta(z,u)=
\Bigl(\prod_{j=1}^{J}(1-u^{E_j}z^{\ell_j})^{-m_j}\Bigr)\,
\zeta_{\mathrm{core}}(z,u),
\]
其中
\[
\ell_j,E_j,m_j\in \ZZ_{>0}\qquad (1\le j\le J),
\]
并设 \(\zeta_{\mathrm{core}}(z,u)\) 的主极点为 \(z_\star(u)\)，且对每个 \(j\) 都有
\[
1-u^{E_j}z_\star(u)^{\ell_j}\neq 0.
\]
再记对应 primitive 生成函数与 Abel 有限部为
\[
P(z,u):=\sum_{n\ge 1}p_n(u)z^n,
\qquad
P_{\mathrm{core}}(z,u):=\sum_{n\ge 1}p_n^{\mathrm{core}}(u)z^n,
\]
\[
\log\mathfrak M(u),
\qquad
\log\mathfrak M_{\mathrm{core}}(u).
\]
则成立：
\begin{enumerate}
  \item \(\zeta(z,u)\) 与 \(\zeta_{\mathrm{core}}(z,u)\) 具有相同的主极点位置与相同的极点阶数；
  \item 若该主极点为单极点，则
  \[
  \operatorname*{Res}_{z=z_\star(u)}\zeta(z,u)
  =
  \operatorname*{Res}_{z=z_\star(u)}\zeta_{\mathrm{core}}(z,u)\,
  \prod_{j=1}^{J}(1-u^{E_j}z_\star(u)^{\ell_j})^{-m_j};
  \]
  \item primitive 生成函数满足
  \[
  P(z,u)=P_{\mathrm{core}}(z,u)+\sum_{j=1}^{J}m_j\,u^{E_j}z^{\ell_j};
  \]
  \item Abel 有限部满足
  \[
  \log\mathfrak M(u)
  =
  \log\mathfrak M_{\mathrm{core}}(u)
  +
  \sum_{j=1}^{J}m_j\,u^{E_j}z_\star(u)^{\ell_j}.
  \]
\end{enumerate}
\end{theorem}
\begin{proof}
记
\[
H(z,u):=\prod_{j=1}^{J}(1-u^{E_j}z^{\ell_j})^{-m_j}.
\]
由假设，
\[
H(z_\star(u),u)\neq 0,
\]
故 \(H(\cdot,u)\) 在 \(z_\star(u)\) 的某邻域内解析且不为零。于是
\[
\zeta(z,u)=H(z,u)\,\zeta_{\mathrm{core}}(z,u)
\]
立即推出：乘子 \(H\) 既不移动主极点，也不改变其阶数，从而第 \(1\) 条成立；若主极点为单极点，则留数只被 \(H(z_\star(u),u)\) 乘上，故第 \(2\) 条成立。

对第 \(3\)、\(4\) 条，只需对每个原子因子反复应用命题 \ref{prop:primitive-orbit-surgery}。每引入一个因子 \((1-u^{E_j}z^{\ell_j})^{-1}\)，primitive 层恰增加一项
\[
u^{E_j}z^{\ell_j},
\]
而 Abel 有限部恰增加一项
\[
u^{E_j}z_\star(u)^{\ell_j}.
\]
重复 \(m_j\) 次便得到系数 \(m_j\)，再对 \(j=1,\dots,J\) 求和，即得
\[
P(z,u)=P_{\mathrm{core}}(z,u)+\sum_{j=1}^{J}m_j\,u^{E_j}z^{\ell_j},
\]
以及
\[
\log\mathfrak M(u)
=
\log\mathfrak M_{\mathrm{core}}(u)
+
\sum_{j=1}^{J}m_j\,u^{E_j}z_\star(u)^{\ell_j}.
\]
证毕。
\end{proof}

\noindent
由此，有限个原子 Euler 因子的手术从不改变主增长率、主极点位置或极点阶数；它们只在单极点情形下以乘法方式重标定主留数，并在 primitive 与 Abel 有限部两层留下严格可加的有限修正。于是，primitive surgery 的作用被压缩为“主部刚性、有限部可加”的统一闭式法则。

\begin{theorem}[有限原子手术在长度加法谐波上的完全对角化]\label{thm:conclusion-finite-primitive-surgery-additive-harmonic-diagonalization}
\leanverified{paper\_conclusion\_finite\_primitive\_surgery\_additive\_harmonic\_diagonalization}
在定理 \ref{thm:conclusion-finite-primitive-surgery-pole-finitepart-transform} 的设定下，对任意整数 \(q\ge 1\) 与任意剩余类 \(a\in \ZZ/q\ZZ\)，定义长度加法投影
\[
\Pi_{a,q}f(z):=\frac{1}{q}\sum_{t=0}^{q-1}\omega_q^{-at}f(\omega_q^t z),
\qquad
\omega_q:=e^{2\pi i/q}.
\]
则 primitive 手术差满足精确对角化公式
\[
\Pi_{a,q}\!\bigl(P-P_{\mathrm{core}}\bigr)(z,u)
=
\sum_{\ell_j\equiv a\ (\mathrm{mod}\ q)}m_j\,u^{E_j}z^{\ell_j}.
\]
更一般地，对任意加法特征 \(\chi(n)=\omega_q^{tn}\) 有
\[
\sum_{n\ge 1}\bigl(p_n(u)-p_n^{\mathrm{core}}(u)\bigr)\chi(n)z^n
=
\sum_{j=1}^{J}m_j\,\chi(\ell_j)\,u^{E_j}z^{\ell_j}.
\]
因此，有限原子手术在长度加法谐波上完全对角化：每个原子只在与其长度同余的唯一 residue channel 中留下质量，而在加法特征下仅留下相位因子 \(\chi(\ell_j)\)。
\end{theorem}
\begin{proof}
由定理 \ref{thm:conclusion-finite-primitive-surgery-pole-finitepart-transform}，
\[
P(z,u)-P_{\mathrm{core}}(z,u)=\sum_{j=1}^{J}m_j\,u^{E_j}z^{\ell_j}.
\]
对两边作用 \(\Pi_{a,q}\)，只需逐项计算单项式的像。对任意 \(\ell\ge 1\)，有
\[
\Pi_{a,q}(z^\ell)
=
\frac{1}{q}\sum_{t=0}^{q-1}\omega_q^{-at}(\omega_q^t z)^\ell
=
\left(\frac{1}{q}\sum_{t=0}^{q-1}\omega_q^{(\ell-a)t}\right)z^\ell.
\]
离散 Fourier 正交关系给出
\[
\frac{1}{q}\sum_{t=0}^{q-1}\omega_q^{(\ell-a)t}
=
\mathbf 1_{\ell\equiv a\ (\mathrm{mod}\ q)},
\]
故
\[
\Pi_{a,q}\!\bigl(P-P_{\mathrm{core}}\bigr)(z,u)
=
\sum_{\ell_j\equiv a\ (\mathrm{mod}\ q)}m_j\,u^{E_j}z^{\ell_j}.
\]
这便得到第一式。

对第二式，只需把
\[
\sum_{n\ge 1}\bigl(p_n(u)-p_n^{\mathrm{core}}(u)\bigr)\chi(n)z^n
\]
看作对 \(\sum_{n\ge 1}(p_n-p_n^{\mathrm{core}})z^n\) 施加长度方向的 Fourier 线性组合。由于 \(\chi(\ell_j)\) 正是单项式 \(z^{\ell_j}\) 在该加法特征下获得的相位，便得
\[
\sum_{n\ge 1}\bigl(p_n(u)-p_n^{\mathrm{core}}(u)\bigr)\chi(n)z^n
=
\sum_{j=1}^{J}m_j\,\chi(\ell_j)\,u^{E_j}z^{\ell_j}.
\]
证毕。
\end{proof}

\begin{corollary}[真实输入 \texorpdfstring{$40$}{40} 状态核的原子 Fourier 指纹]\label{cor:conclusion-realinput40-atomic-fourier-fingerprint}
在真实输入 \(40\) 状态核中，
\[
\Delta(z,u)=(1-uz^2)F(z,u),
\qquad
P(z,u)=uz^2+P_{\mathrm{core}}(z,u).
\]
记 \(z_\star(u)\) 为由 core 极点 \(F(z_\star(u),u)=0\) 所确定的主极点，并令 \(b\) 为方程
\[
1-6b+9b^2-b^3-b^4=0
\]
的最大正根。对任意有限加法特征 \(\chi:\ZZ/q\ZZ\to \CC^\times\)，定义原子 Fourier 振幅
\[
\mathcal A_\chi(u)
:=
\left[\sum_{n\ge 1}\bigl(p_n(u)-p_n^{\mathrm{core}}(u)\bigr)\chi(n)z^n\right]_{z=z_\star(u)}.
\]
则有精确公式
\[
\mathcal A_\chi(u)=\chi(2)\,u\,z_\star(u)^2,
\]
并在低温端 \(u\to +\infty\) 满足
\[
\mathcal A_\chi(u)\longrightarrow b\,\chi(2).
\]
等价地，对 residue-filtered 振幅
\[
\mathcal A_{a,q}(u):=
\Pi_{a,q}\!\bigl(P-P_{\mathrm{core}}\bigr)(z_\star(u),u),
\]
有
\[
\mathcal A_{a,q}(u)=u\,z_\star(u)^2\,\mathbf 1_{a\equiv 2\ (\mathrm{mod}\ q)},
\qquad
\mathcal A_{a,q}(u)\longrightarrow b\,\mathbf 1_{a\equiv 2\ (\mathrm{mod}\ q)}.
\]
因此，真实输入模型中的唯一原子在零温极限中留下一个秩一 Fourier 指纹：全部长度谐波仅在频率 \(2\) 处存活，且极限质量恰为 \(b\)。
\end{corollary}
\begin{proof}
将定理 \ref{thm:conclusion-finite-primitive-surgery-additive-harmonic-diagonalization} 应用于唯一原子
\[
(m,\ell,E)=(1,2,1),
\]
便得到
\[
\sum_{n\ge 1}\bigl(p_n(u)-p_n^{\mathrm{core}}(u)\bigr)\chi(n)z^n
=
\chi(2)\,u z^2.
\]
在 \(z=z_\star(u)\) 处取值，即得
\[
\mathcal A_\chi(u)=\chi(2)\,u\,z_\star(u)^2.
\]
而本文前面已证明
\[
u\,z_\star(u)^2\longrightarrow b
\qquad
(u\to +\infty),
\]
故
\[
\mathcal A_\chi(u)\longrightarrow b\,\chi(2).
\]

对 residue-filtered 振幅，只需把 \(\chi\) 换成投影 \(\Pi_{a,q}\) 即可。由第一条结论立刻得到
\[
\mathcal A_{a,q}(u)
=
u\,z_\star(u)^2\,\mathbf 1_{a\equiv 2\ (\mathrm{mod}\ q)},
\]
并在 \(u\to+\infty\) 时收敛到
\[
b\,\mathbf 1_{a\equiv 2\ (\mathrm{mod}\ q)}.
\]
证毕。
\end{proof}

\begin{theorem}[relative \texorpdfstring{$\zeta$}{zeta} 商的临界阶不变性]\label{thm:conclusion-relative-zeta-quotient-critical-order-invariance}
\leanverified{paper\_conclusion\_relative\_zeta\_quotient\_critical\_order\_invariance}
设 \(B\) 为目标核，\(M\) 为参考核，且 \(\rho(M)<3\)。定义相对对象
\[
\zeta_{B/M}(z):=\frac{\zeta_B(z)}{\zeta_M(z)},
\qquad
r_n:=a_n(B)-a_n(M),
\qquad
q_n:=p_n(B)-p_n(M).
\]
若 \(\zeta_B\) 在 \(z_\ast=\frac13\) 处具有简单极点，而 \(\zeta_M\) 在 \(z_\ast\) 解析且 \(\zeta_M(z_\ast)\neq 0\)，则
\[
\zeta_{B/M}(z)=\frac{C_{B/M}}{1-3z}+O(1),
\qquad
C_{B/M}=\frac{C_B}{\zeta_M(1/3)}.
\]
因此，相对化只能重整化留数与有限部，不能改变临界极点阶数。
\end{theorem}

\begin{proof}
由于 \(\rho(M)<3\)，参考核 \(\zeta_M\) 在 \(z_\ast=1/3\) 的邻域内解析；再加上 \(\zeta_M(z_\ast)\neq 0\)，其在该点局部可逆。把 \(\zeta_B\) 的简单极点展开
\[
\zeta_B(z)=\frac{C_B}{1-3z}+O(1)
\]
除以解析且非零的 \(\zeta_M(z)\)，便得到
\[
\zeta_{B/M}(z)=\frac{C_B}{\zeta_M(1/3)}\frac{1}{1-3z}+O(1).
\]
故 relative \(\zeta\) 商与目标核共享同一简单极点阶，只在留数上被 \(\zeta_M(1/3)^{-1}\) 缩放。证毕。
\end{proof}

\begin{corollary}[参考减法无法抹除本原相的主指数]\label{cor:conclusion-relative-primitive-subtraction-preserves-main-exponent}
\leanverified{paper\_conclusion\_relative\_primitive\_subtraction\_preserves\_main\_exponent}
在定理 \ref{thm:conclusion-relative-zeta-quotient-critical-order-invariance} 的设定下，若 \(C_{B/M}\neq 0\)，则相对 primitive 差异
\[
q_n=p_n(B)-p_n(M)
\]
与目标 primitive 计数 \(p_n(B)\) 共享同一临界指数尺度。换言之，参考核只能改变振幅与 Abel 有限部，不能消去主指数。
\end{corollary}

\begin{proof}
定理 \ref{thm:conclusion-relative-zeta-quotient-critical-order-invariance} 已证明：\(\zeta_{B/M}\) 在 \(z_\ast=1/3\) 仍保留同阶简单极点。按本文 primitive/\(\zeta\)/Abel 接口的标准传递，简单极点阶决定 primitive 计数的主指数尺度，而留数只控制主振幅及有限部常数的位移。故只要 \(C_{B/M}\neq 0\)，序列 \((q_n)\) 与 \((p_n(B))\) 具有同一主指数。证毕。
\end{proof}

\begin{theorem}[真实输入 \texorpdfstring{$40$}{40} 状态核的 prime spike--background 律]\label{thm:conclusion-realinput40-prime-spike-background-law}
在真实输入 \(40\) 状态输出位势核中，
\[
P(z,u)=P_{\mathrm{core}}(z,u)+u z^2.
\]
沿用定理 \ref{thm:conclusion-ramanujan-finite-cyclotomic-shell-projection} 的记号，则对任意 \(Q\ge 1\) 都有
\[
\mathcal R_{Q,r}(z)
=
\mathcal R_{Q,r}^{\mathrm{core}}(z)
\;+\;
c_Q(1-r)\,z^2.
\]
特别地，当 \(Q=p\) 为素数时，
\[
\mathcal R_{p,r}(z)
=
\mathcal R_{p,r}^{\mathrm{core}}(z)
\;+\;
\begin{cases}
(p-1)z^2, & r\equiv 1\pmod p,\\[3pt]
-z^2, & r\not\equiv 1\pmod p.
\end{cases}
\]
并且
\[
\sum_{r\ (\mathrm{mod}\ p)}
\bigl(\mathcal R_{p,r}(z)-\mathcal R_{p,r}^{\mathrm{core}}(z)\bigr)=0.
\]
\leanverified{paper\_conclusion\_realinput40\_prime\_spike\_background\_law}
\end{theorem}
\begin{proof}
由真实输入 \(40\) 状态核的 primitive 手术公式，
\[
P(z,u)-P_{\mathrm{core}}(z,u)=u z^2,
\]
这正对应唯一原子
\[
(m,\ell,E)=(1,2,1).
\]
将定理 \ref{thm:conclusion-ramanujan-finite-cyclotomic-shell-projection} 应用于该唯一原子，即得
\[
\mathcal R_{Q,r}(z)-\mathcal R_{Q,r}^{\mathrm{core}}(z)=c_Q(1-r)\,z^2.
\]
若 \(Q=p\) 为素数，则
\[
c_p(0)=p-1,
\qquad
c_p(n)=-1\quad (p\nmid n),
\]
故得到 spike--background 公式。最后
\[
\sum_{r\ (\mathrm{mod}\ p)}c_p(1-r)=0
\]
给出总和守恒。证毕。
\end{proof}

\begin{proposition}[七条几何两步回路向单一 Ramanujan 单项式的严格塌缩]\label{prop:conclusion-realinput40-seven-two-step-loops-collapse}
在同一 \(40\) 状态真实输入核中，branch 审计给出的满足
\[
(|\tau|,E)=(2,1)
\]
的几何两步回路共计 \(7\) 条；然而在 primitive--Witt 层，它们的净效应严格塌缩为单一 monomial
\[
u z^2.
\]
因而经 Ramanujan 投影后，其全部算术遗迹仅为
\[
c_Q(1-r)\,z^2.
\]
\end{proposition}
\begin{proof}
文稿中的 branch 审计已给出：分支层中 \((|\tau|,E)=(2,1)\) 的几何两步回路共有 \(7\) 条。另一方面，真实输入模型的 primitive 生成函数满足
\[
P(z,u)-P_{\mathrm{core}}(z,u)=u z^2,
\]
并非 \(7u z^2\)。这说明几何七重性在 primitive--Witt 商中完全塌缩为一个长度--能量原子坐标。

再由定理 \ref{thm:conclusion-realinput40-prime-spike-background-law}，该唯一原子的 Ramanujan 投影恰是
\[
c_Q(1-r)\,z^2.
\]
证毕。
\end{proof}

\begin{theorem}[有限 cyclotomic 壳层的双重律：指数静默，有限部活跃]\label{thm:conclusion-finite-cyclotomic-shell-double-law}
设
\[
R_u(z):=\sum_{j=1}^{N}m_j\,u^{E_j}z^{\ell_j},
\qquad
P(z,u)=P_{\mathrm{core}}(z,u)+R_u(z).
\]
则对每个固定的 \(u\) 都成立：
\begin{enumerate}
  \item 若
  \[
  L:=\max_{1\le j\le N}\ell_j,
  \]
  则
  \[
  p_n(u)=p_n^{\mathrm{core}}(u)\qquad (n>L).
  \]
  因而
  \[
  \limsup_{n\to\infty}|p_n(u)|^{1/n}
  =
  \limsup_{n\to\infty}|p_n^{\mathrm{core}}(u)|^{1/n},
  \]
  并且一切仅由 tail-exponential 行为决定的不变量，包括收敛半径、dominant pole exponent 与 RH/Ramanujan 型平方根门槛，全部壳层不变。
  \item 若主增长率 \(\lambda(u)\) 由 core 极点决定，则 Abel 有限部常数满足精确位移
  \[
  \log\mathfrak M(u)-\log\mathfrak M_{\mathrm{core}}(u)
  =
  \sum_{j=1}^{N}m_j\,\frac{u^{E_j}}{\lambda(u)^{\ell_j}}.
  \]
\end{enumerate}
特别地，对真实输入 \(40\) 状态核，
\[
\log\mathfrak M(u)-\log\mathfrak M_{\mathrm{core}}(u)=u\,z_\star(u)^2,
\]
在 \(u=1\) 处回收到文稿审计得到的位移 \(\varphi^{-4}\)。
\end{theorem}
\begin{proof}
第 \(1\) 条直接来自
\[
P(z,u)-P_{\mathrm{core}}(z,u)=R_u(z),
\]
其中 \(R_u(z)\) 是一个有限多项式。故只有有限多个 primitive 系数发生改变，从而所有 tail-exponential 不变量均与 core 完全相同。

第 \(2\) 条正是定理 \ref{thm:conclusion-finite-primitive-surgery-pole-finitepart-transform} 的 Abel 有限部变换公式在
\[
z_\star(u)=\lambda(u)^{-1}
\]
处的改写。

对真实输入 \(40\) 状态核，唯一原子为 \((m,\ell,E)=(1,2,1)\)，故
\[
\log\mathfrak M(u)-\log\mathfrak M_{\mathrm{core}}(u)=u\,z_\star(u)^2.
\]
在 \(u=1\) 时，文稿已审计出 \(z_\star(1)=\varphi^{-2}\)，于是位移即为 \(\varphi^{-4}\)。证毕。
\end{proof}

\noindent
因此，有限 cyclotomic 壳层在指数律上完全静默，却对 Abel 有限部留下严格可加、可审计、可分层的位移。对真实输入 \(40\) 状态核而言，这条双重律又与 prime spike--background 公式闭合起来：几何七重性的全部算术遗迹最终只剩下一枚 \(c_Q(1-r)z^2\) 型净荷。


\endinput
